\documentclass[conference]{IEEEtran}
\IEEEoverridecommandlockouts

% Core mathematics and symbols
\usepackage{amsmath,amssymb,amsfonts}
\usepackage{textcomp}

% Graphics and tables
\usepackage[dvipsnames,svgnames,table]{xcolor}
\usepackage{graphicx}
\usepackage{array}
\newcolumntype{H}{>{\setbox0=\hbox\bgroup}c<{\egroup}@{}}
\usepackage{booktabs}
\usepackage{siunitx}
\DeclareSIUnit{\dBm}{dBm}
\DeclareSIUnit{\dBi}{dBi}
\DeclareSIUnit{\dBsm}{dBsm}
\sisetup{detect-all}

% Bibliography and cross references
\usepackage[backend=biber,language=english,autolang=other,style=ieee]{biblatex}
\addbibresource{bib.bib}
\AtBeginBibliography{\footnotesize}
\usepackage[hidelinks]{hyperref}
\usepackage[capitalize]{cleveref}

% Plotting support
\usepackage{pgfplots}
\pgfplotsset{compat=1.18}
\usepgfplotslibrary{groupplots,dateplot,statistics}
\usepackage{pgfplotstable}

% Glossaries and author comments
\usepackage[xindy,nonumberlist]{glossaries}
\makenoidxglossaries
\loadglsentries{abbr}
\input{auto_names}
\usepackage{orcidlink}
\usepackage[font=small]{caption}
\usepackage{subcaption}
\usepackage{placeins}

% IEEE specific
\def\BibTeX{{\rm B\kern-.05em{\sc i\kern-.025em b}\kern-.08em
    T\kern-.1667em\lower.7ex\hbox{E}\kern-.125emX}}


%%%%%%%%%%%%%%%%%%%%%%%%%%%% COLORS FOR TIKZ %%%%%%%%%%%%%%%%%%%%%%%%%%%%%%%%%%%%
\definecolor{color1}{HTML}{1b9e77}
\definecolor{color2}{HTML}{d95f02}
\definecolor{color3}{HTML}{7570b3}
\definecolor{color4}{HTML}{e7298a}
\definecolor{color5}{HTML}{66a61e}
\definecolor{color6}{HTML}{e6ab02}
\definecolor{color7}{HTML}{a6761d}
\definecolor{color8}{HTML}{666666}
%%%%%%%%%%%%%%%%%%%%%%%%%%%% COLORS FOR TIKZ %%%%%%%%%%%%%%%%%%%%%%%%%%%%%%%%%%%%


%%%%%%%%%%%%%% DEFINE YOUR OWN COLORS HERE %%%%%%%%%%%%%%%%%%%
% you can use \colorlet to make synonyms, so you can use your own naming conventions
\colorlet{AMBgreen}{color1}
\colorlet{AMBdarkgreen}{color1}
\colorlet{AMBlightgreen}{color1}
\definecolor{darkgray176}{RGB}{176,176,176}

% Strategy colors and compact labels used in PGFPlots figures.
\definecolor{CentralMassiveMimoColor}{HTML}{D4A017}
\definecolor{DistributedFixedColor}{HTML}{2F5D8A}
\definecolor{DistributedMovableColor}{HTML}{CB3A2A}
\definecolor{DistributedMovableOptTwoColor}{HTML}{1B9E77}
\definecolor{DistributedMovableOptThreeColor}{HTML}{7570B3}




\input{definitions/def-tikz}
\newcommand{\resultroot}{../outputs/rabot}
\newcommand{\postroot}{\resultroot/postprocessing}

\newcommand{\ConceptFigurePath}{figures/cocoon_concept.pdf}
\newcommand{\SceneLayoutFigurePath}{\resultroot/scene_layout.png}
\newcommand{\TrajectoryColormapFigurePath}{\resultroot/trajectory_colormap.png}
\newcommand{\SINRCdfFigurePath}{\postroot/sinr/sinr_cdf_publication.png}
\newcommand{\SINRWorstUserFigurePath}{\postroot/sinr/worst_user_sinr_timeseries.png}
\newcommand{\SINRUsersBelowFigurePath}{\postroot/sinr/users_below_threshold_timeseries.png}
\newcommand{\SINRThresholdSweepFigurePath}{\postroot/sinr/sinr_threshold_sweep.png}
\newcommand{\ScheduleOverviewFigurePath}{\postroot/schedule/schedule_strategy_overview.png}
\newcommand{\ScheduleHistogramFigurePath}{\postroot/schedule/schedule_transition_distance_histogram.png}
\newcommand{\CoverageFigurePath}{\resultroot/coverage_map.png}
\newcommand{\FixedCoverageFigurePath}{\resultroot/fixed_coverage_map.png}

\NewDocumentCommand{\MissingFigureBox}{ O{0.95\linewidth} m }{%
  \setlength{\fboxsep}{0.8em}%
  \fbox{%
    \parbox[c][0.22\textheight][c]{#1}{%
      \centering
      \footnotesize
      Missing figure asset\\[0.4em]
      \texttt{\detokenize{#2}}\\[0.4em]
      Rebuild or copy the expected output before finalizing the manuscript.
    }%
  }%
}

\NewDocumentCommand{\PaperImage}{ O{\linewidth} m }{%
  \IfFileExists{#2}{\includegraphics[width=#1]{#2}}{\MissingFigureBox[#1]{#2}}%
}

\NewDocumentCommand{\PaperFigure}{ O{tb} O{0.95\linewidth} m m m }{%
  \begin{figure}[#1]
    \centering
    \PaperImage[#2]{#3}
    \caption{#4}
    \label{#5}
  \end{figure}%
}

\NewDocumentCommand{\PaperSubfigurePair}{ O{tb} O{0.48\textwidth} m m m m m m m m }{%
  \begin{figure*}[#1]
    \centering
    \begin{subfigure}[t]{#2}
      \centering
      \PaperImage[\linewidth]{#3}
      \caption{#4}
      \label{#5}
    \end{subfigure}\hfill
    \begin{subfigure}[t]{#2}
      \centering
      \PaperImage[\linewidth]{#6}
      \caption{#7}
      \label{#8}
    \end{subfigure}
    \caption{#9}
    \label{#10}
  \end{figure*}%
}



\defineauthors[false]{gilles}

\begin{document}

\title{Towards Cocoon: A Simulation Scaffold for Citizen-Co-Created Low-Height Distributed MIMO Access Networks%
\thanks{Working manuscript scaffold for the COCOON project. Update title, funding text, and author list before submission.}%
\thanks{For the purpose of open access, the author has applied a CC BY public copyright license to any Author Accepted Manuscript version arising from this submission.}%
}

\author{
\IEEEauthorblockN{%
Gilles Callebaut\,\orcidlink{0000-0003-2413-986X}\IEEEauthorrefmark{1},
Other Author\,\orcidlink{0000-0003-2413-986X}\IEEEauthorrefmark{2},
}
\IEEEauthorblockA{\IEEEauthorrefmark{1}
KU Leuven, Belgium}
\IEEEauthorblockA{\IEEEauthorrefmark{2}
University, Country}
}

\maketitle

\begin{abstract}
This manuscript is structured as a concept-plus-simulation paper for the Cocoon vision. The current draft focuses on the simulation scaffold built in \texttt{COCOON-Sionna}: an outdoor \gls{dmimo} pipeline that generates low-height candidate \glspl{ap}, simulates user mobility, computes path-based \gls{csi}, and compares movable-\gls{ap} placement strategies over a common urban scenario. The interpretation of the generated metrics is intentionally deferred. Instead, this draft provides the paper structure, figure hooks, and the core technical framing needed to start incorporating results from committed simulation outputs.\gilles{Replace this abstract with quantitative claims once the first stable result set is frozen.}
\end{abstract}

\begin{IEEEkeywords}
cell-free systems, distributed MIMO, low-height access points, access-point placement, urban mobility, Sionna RT
\end{IEEEkeywords}

\glsresetall

\section{Introduction}

The Cocoon project targets a wireless access paradigm in which the network adapts to the fluidity of the city rather than assuming a static infrastructure and a stationary demand profile. The motivating hypothesis from the project proposal is that a practical \gls{cf} network becomes more feasible when low-height infrastructure can be co-created, repositioned, and activated only when the measured channel conditions justify it. In that view, the citizens are not only traffic sources but also anchors for sensing, mounting, and reconfiguring the surrounding network.

This article narrows that broader vision to a simulation-based study of low-height infrastructure adaptation. The implemented framework covers wall-mounted candidate \glspl{ap}, outdoor user mobility, path-based \gls{csi} extraction, and repeated placement decisions over relocation windows in an urban scene. This scope enables a reproducible comparison of geometry, radio-performance, and relocation-cost metrics for alternative infrastructure strategies.

\PaperFigure[t][0.86\linewidth]{\ConceptFigurePath}{Conceptual illustration of the Cocoon low-height access-network vision, where citizens interact with the network and move \glspl{ap} depending on the current channel conditions and its users. By dynamically adapting to the current network state and requirement, the ecological best solution can be achieved when all network resources can cooperate in terms of energy efficiency and minimum materials needed. The lines between the people represent the real-time channel gain measurements. Each person thereby represents a potential \gls{ap} candidate position. Based on these measurements, the most optimal \gls{ap} location can be determined, which is in this case the person close to the tree.}{fig:cocoon_concept}

The main contributions of the present study are:
\begin{itemize}
    \item a concise technical framing of Cocoon as a low-height, movable-\gls{ap}, citizen-co-created access-network concept,
    \item a reproducible description of the Rabot outdoor simulation pipeline implemented in \texttt{COCOON-Sionna}, and
    \item a common analysis workflow for geometry, \gls{sinr}, \gls{esr}, and relocation-cost results derived from exported simulation artifacts.
\end{itemize}

\section{Scope Within Cocoon}

The full Cocoon proposal spans citizen interaction with the wireless network, channel assessment, architecture design, and proof-of-concept validation. The present article addresses the simulation-facing part of that program, with emphasis on dynamic channel-aware adaptation, low-height \gls{ap}/\gls{csp} placement, and the algorithmic consequences of allowing physical repositioning.

The main modeling simplification is that co-creation is represented here through an abstract movable infrastructure layer rather than through a user study or a deployment trial. Candidate mounting locations stand in for positions that could later be offered to citizens, and the placement scheduler stands in for the future tool chain that would notify participants where the next best \gls{csp} constellation should be installed.

Accordingly, the scope of the present article is:
\begin{itemize}
    \item the paper validates a simulation methodology and a comparison workflow, not the social feasibility of citizen participation,
    \item the paper studies a Rabot outdoor scenario, not yet a multi-site or multi-season campaign,
    \item the paper compares placement strategies over a common \gls{csi}-derived objective, not yet a full end-to-end network stack.
\end{itemize}

\subsection{Placement Strategies}
\label{sec:placement_strategies}

The placement layer compares the following strategy definitions under a common \gls{ap} budget:
\begin{itemize}
    \item[] \strategy1 -- one rooftop proxy is placed on the building whose rooftop representative point is nearest the area center and is oriented toward that center. Its antenna topology is the normalized co-located $\num{6}\times\num{6}$ \gls{upa} obtained by aggregating the distributed element budget into one rooftop site; it is used only as the central reference.
    \item[] \strategy2 -- one distributed subset is sampled from the wall-mounted candidate pool and kept unchanged over the full trajectory.
    \item[] \strategy4 -- the distributed subset is reselected every relocation window from a nearest-user \gls{ue}-to-\gls{ue} proxy-\gls{csi} model. Each candidate site binds to the closest user within $d_{\mathrm{th}}$, the controller forms a wideband proxy channel from that user's measured \gls{ue}-to-\gls{ue} CFR, and the subset is ranked by proxy \gls{esr}.
    \item[] \strategy5 -- the distributed subset is reselected every relocation window from the same nearest-user distance-threshold proxy model, but the controller uses the stored \gls{ue}-to-\gls{ue} link-power row of the anchor user, replaces the diagonal entry by the strongest non-self value, rescales by $P_{\mathrm{AP}}/P_{\mathrm{UE}}$, and ranks the subset by proxy average power.
\end{itemize}

\section{System Model}
\label{sec:system_model}

We consider a downlink \gls{dmimo} system with active \glspl{ap} $\mathcal{A}=\{1,\ldots,M\}$ and scheduled \glspl{ue} $\mathcal{U}=\{1,\ldots,K\}$. Let $N_m$ denote the number of transmit ports at \gls{ap} $m$ and $N_{r,k}$ the number of receive ports at \gls{ue} $k$. The per-\gls{ap} channel block is
\begin{equation}
    \mathbf{H}_{k,m} \in \mathbb{C}^{N_{r,k} \times N_m},
\end{equation}
and the full distributed channel observed by \gls{ue} $k$ is
\begin{equation}
    \mathbf{H}_k = \left[\mathbf{H}_{k,1}\;\mathbf{H}_{k,2}\;\cdots\;\mathbf{H}_{k,M}\right].
\end{equation}

Given a receive combiner $\mathbf{u}_k$, the effective channel becomes
\begin{equation}
    \mathbf{h}_k^{\mathsf{H}} = \mathbf{u}_k^{\mathsf{H}}\mathbf{H}_k.
\end{equation}
Stacking all users yields the effective distributed channel matrix $\mathbf{H}_{\mathrm{eff}}$. The current pipeline applies a least-squares pseudo-inverse \gls{zf} precoder on $\mathbf{H}_{\mathrm{eff}}$, followed by per-stream power allocation under a total \gls{ap} power budget.

For \gls{ue} $k$, the instantaneous downlink \gls{sinr} is evaluated as
\begin{equation}
    \mathrm{SINR}^{\mathrm{DL}}_k =
    \frac{p_k \left|\mathbf{h}_k^{\mathsf{H}}\mathbf{f}_k\right|^2}
    {\sum_{i \in \mathcal{U},\, i \neq k} p_i \left|\mathbf{h}_k^{\mathsf{H}}\mathbf{f}_i\right|^2 + \sigma^2},
\end{equation}
where $\mathbf{f}_i$ is the precoding vector for stream $i$, $p_i$ its allocated power, and $\sigma^2$ the thermal-noise term. The implementation keeps the interference term explicitly and exports desired-signal, interference, and noise power terms for each evaluated user snapshot.

The \gls{esr} used later in the results section is derived from the same per-snapshot \gls{sinr} values. For snapshot $n$ with scheduled-user set $\mathcal{U}_n$, the postprocessing layer computes
\begin{equation}
    \mathrm{ESR}_n = \sum_{k \in \mathcal{U}_n}\log_2\!\left(1+\mathrm{SINR}^{\mathrm{DL}}_{k,n}\right).
\end{equation}
These per-snapshot \gls{sinr} values form the basis of the trajectory-level \gls{sinr} and \gls{esr} statistics reported later in the article.

\subsection{Placement Strategies}

The placement layer compares five strategies under a common \gls{ap} budget:
\begin{itemize}
    \item \texttt{central\_massive\_mimo}: one rooftop proxy \gls{ap} is placed on the building whose rooftop representative point is nearest the area center; the proxy is oriented toward the area center. This is a normalized central reference, not a full single co-located-array placement model.
    \item \texttt{distributed\_fixed}: a distributed set of wall-mounted candidate positions is sampled once and kept fixed over the full trajectory.
    \item \texttt{distributed\_movable}: optimization 1. The same number of distributed \glspl{ap} is reselected at each relocation window from the wall-mounted candidate pool with a history-aware local-\gls{csi} rule.
    \item \texttt{distributed\_movable\_optimization\_2}: optimization 2. The same number of distributed \glspl{ap} is reselected at each relocation window from the wall-mounted candidate pool with a current-window nearest-user \gls{ue}-to-\gls{ue} proxy-\gls{csi} rule and a proxy-\gls{esr} objective.
    \item \texttt{distributed\_movable\_optimization\_3}: optimization 3. The same number of distributed \glspl{ap} is reselected at each relocation window from the wall-mounted candidate pool with the same nearest-user proxy association as optimization 2, but with a proxy-average-power objective based on $|\mathrm{CSI}|^2$.
\end{itemize}

Let $\mathcal{C}$ denote the wall-mounted candidate pool and let $\mathcal{S}_w \subseteq \mathcal{C}$ denote the selected movable-\gls{ap} subset in relocation window $w$. In the Rabot scenario studied here, $\lvert \mathcal{S}_w \rvert = N_{\mathrm{mov}}=6$ and no always-on distributed \glspl{ap} are used, so the distributed comparison isolates the effect of changing the movable subset over time.

\subsection{Relocation Heuristic}

At an intuitive level, the relocation mechanism repeats the same observe-score-relocate loop at every relocation instant. Optimization~1 uses simulated \gls{ap}-to-\gls{ue} channel observations from previously evaluated subsets and discounts older samples exponentially. Optimization~2 and optimization~3 follow a different practical assumption: the \gls{ap}-to-\gls{ue} channel of a dormant candidate position is not known before an \gls{ap} is actually moved there. Consequently, these two rules do not optimize directly over candidate-position \gls{ap}-to-\gls{ue} \gls{csi}. Instead, each candidate site is associated with the closest user within a distance threshold, the measured \gls{ue}-to-\gls{ue} channel of that nearby user is reused as a proxy for the dormant candidate site, and the subset is ranked either by proxy \gls{esr} (optimization~2) or by proxy average received power (optimization~3). Once the best subset has been chosen, the movable \glspl{ap} are relocated to those candidate positions and the actual deployed \gls{ap}-to-\gls{ue} \gls{csi} is recomputed for storage, postprocessing, and final reporting.

The movable distributed strategy is implemented as a windowed receding-horizon controller. Let $\{t_n\}_{n=0}^{N_t-1}$ denote the simulated snapshot times and let $T_{\mathrm{rel}}$ denote the relocation interval. The trajectory is partitioned into disjoint decision windows
\begin{equation}
    \mathcal{W}_w = \left\{n \in \{0,\ldots,N_t-1\} :
    \left\lfloor \frac{t_n-t_0}{T_{\mathrm{rel}}}\right\rfloor = w\right\}.
\end{equation}
At the beginning of each window $w$, the controller selects a subset $\mathcal{S}_w \subseteq \mathcal{C}$ with cardinality $\lvert \mathcal{S}_w \rvert = N_{\mathrm{mov}}$, where $N_{\mathrm{mov}}$ is the movable-\gls{ap} budget. The fixed distributed mode uses one time-invariant subset, whereas the movable mode recomputes $\mathcal{S}_w$ at every relocation instant.

\paragraph{Optimization 1: historical $K$-nearest local-\gls{csi} controller.}
The first online decision rule is intentionally local and does not optimize a full radio map. For a candidate site $c \in \mathcal{C}$ with horizontal position $\mathbf{x}_c \in \mathbb{R}^2$, define the local neighborhood $\mathcal{N}_K(c,w)$ as the set of the $K$ nearest \gls{ue}-snapshot pairs in the horizontal plane within window $w$,
\begin{equation}
    \mathcal{N}_K(c,w) =
    \operatorname{KNN}_K\!\left(
        \mathbf{x}_c;
        \{\mathbf{r}_{u,n} : u \in \mathcal{U}, n \in \mathcal{W}_w\}
    \right),
\end{equation}
where $\mathbf{r}_{u,n}$ denotes the horizontal position of \gls{ue} $u$ at time index $n$. For a subset $\mathcal{S} \subseteq \mathcal{C}$, the support region used by the controller is the union
\begin{equation}
    \mathcal{N}_K(\mathcal{S},w) = \bigcup_{c \in \mathcal{S}} \mathcal{N}_K(c,w).
\end{equation}
In the Rabot configuration, $K=10$, so each candidate contributes at most ten local \gls{ue}-snapshot pairs per window. This restriction makes the relocation rule sensitive to the \gls{csi} observed near the candidate sites rather than to globally averaged scene statistics. In other words, user positions are used to construct the local sampling mask, but they do not enter the score directly.

For every queried subset $\mathcal{S}$, the implementation computes the corresponding \gls{ap}-to-\gls{ue} \gls{csi}/\gls{sinr} segment over the current window and stores it in a per-subset cache. Thus, the relocation heuristic is driven by channel-response information associated with the candidate subset under evaluation, not by geometry alone. If $\gamma_{u,n}(\mathcal{S})$ denotes the best-link \gls{sinr} sample produced when subset $\mathcal{S}$ is active, the historical record available when deciding window $w$ is
\begin{equation}
    \mathcal{H}_w(\mathcal{S}) =
    \left\{(t_n,\gamma_{u,n}(\mathcal{S})) :
    u \in \mathcal{U}, n \in \mathcal{W}_j, j = 0,\ldots,w\right\}.
\end{equation}
When the same subset is revisited later, previously computed window segments are reused and only the new current-window segment is appended. Thus, the historical objective remains consistent while avoiding repeated ray-tracing evaluations for past windows. In the present implementation, the relocation score is formed from these \gls{ap}-to-\gls{ue} channel samples only; \gls{ue}-to-\gls{ue} \gls{csi} is reserved for the final trajectory-level comparison score, and \gls{ap}-to-\gls{ap} \gls{csi} is exported for analysis rather than for online relocation control.

Historical CSI is discounted through exponential aging. If a sample was observed at time $t_n$ and the relocation decision for window $w$ is taken at $t_w^{\mathrm{end}}$, the associated weight is
\begin{equation}
    \omega_n^{(w)} = \exp\!\left[-\lambda\left(t_w^{\mathrm{end}}-t_n\right)\right],
\end{equation}
with $\lambda$ the configured decay rate. Hence, newer samples always dominate older ones: a sample measured at the decision time keeps weight $1$, a sample that is one relocation window old keeps weight $\exp(-\lambda T_{\mathrm{rel}})$, and the weight keeps halving every additional relocation window when $\lambda=\ln(2)/T_{\mathrm{rel}}$, as in the Rabot configuration. The local control objective is then the weighted empirical 10th percentile of the \gls{sinr} samples retained by the local neighborhoods across all observed windows,
\begin{equation}
    J_w^{(1)}(\mathcal{S}) =
    Q_{0.1}^{\omega}\!\left(
        \left\{\gamma_{u,n}(\mathcal{S}) :
        (u,n) \in \bigcup_{j=0}^{w} \mathcal{N}_K(\mathcal{S},j)\right\}
    \right),
\end{equation}
where $Q_{0.1}^{\omega}(\cdot)$ denotes the weighted 10th-percentile operator induced by $\omega_n^{(w)}$. The use of a lower-tail percentile, rather than a mean, makes the controller deliberately conservative with respect to persistently weak local links.

\paragraph{Optimization 2: current-window nearest-user proxy-\gls{csi} controller.}
The second online rule removes the $K$-nearest construction and does not use historical samples. Let $d_{\mathrm{th}}$ denote the configured distance threshold. For each candidate site $c \in \mathcal{C}$ and snapshot $n \in \mathcal{W}_w$, define the anchor user
\begin{equation}
    \ell(c,n)=
    \arg\min_{u \in \mathcal{U}}
    \left\|\mathbf{r}_{u,n}-\mathbf{x}_c\right\|_2,
\end{equation}
and the validity indicator
\begin{equation}
    \delta(c,n)=
    \mathbf{1}\!\left\{
        \left\|\mathbf{r}_{\ell(c,n),n}-\mathbf{x}_c\right\|_2 \le d_{\mathrm{th}}
    \right\}.
\end{equation}
The dormant candidate-position \gls{ap}-to-\gls{ue} channel is not available before relocation. Therefore, optimization~2 replaces it by a proxy channel extracted from the measured \gls{ue}-to-\gls{ue} \gls{csi}. If $h^{\mathrm{UU}}_{u,v,n}$ denotes the narrowband \gls{ue}-to-\gls{ue} coefficient from transmitter-user $v$ to receiver-user $u$ at snapshot $n$, the proxy coefficient used for candidate $c$ is
\begin{equation}
    \tilde{h}_{u,c,n}=
    \begin{cases}
        h^{\mathrm{UU}}_{u,\ell(c,n),n}, & \delta(c,n)=1,\; u \neq \ell(c,n), \\
        h^{\mathrm{UU}}_{v^\star,\ell(c,n),n}, & \delta(c,n)=1,\; u = \ell(c,n), \\
        0, & \delta(c,n)=0,
    \end{cases}
\end{equation}
where $v^\star=\arg\max_{v \in \mathcal{U}} |h^{\mathrm{UU}}_{v,\ell(c,n),n}|$ supplies a surrogate self-link because the diagonal user-to-itself measurement is not observable. For a candidate subset $\mathcal{S}$, the proxy channel matrix $\tilde{\mathbf{H}}_n(\mathcal{S})$ is formed by stacking the columns $\tilde{h}_{u,c,n}$ for all $c \in \mathcal{S}$. The same \gls{zf} procedure as in the deployed \gls{ap} case is then applied to $\tilde{\mathbf{H}}_n(\mathcal{S})$, yielding a proxy \gls{sinr} $\tilde{\gamma}_{u,n}(\mathcal{S})$ and the current-window proxy-\gls{esr} objective
\begin{equation}
    J_w^{(2)}(\mathcal{S}) =
    \frac{1}{|\mathcal{W}_w|}
    \sum_{n \in \mathcal{W}_w}
    \sum_{u \in \mathcal{U}}
    \log_2\!\left(1+\tilde{\gamma}_{u,n}(\mathcal{S})\right).
\end{equation}
Consequently, optimization~2 uses measured \gls{ue}-to-\gls{ue} proxy \gls{csi} to rank candidate subsets, while avoiding the unrealistic assumption that dormant candidate-position \gls{ap}-to-\gls{ue} channels are already known.

\paragraph{Optimization 3: current-window nearest-user proxy-power controller.}
The third online rule reuses the same nearest-user anchor $\ell(c,n)$ and validity indicator $\delta(c,n)$ as optimization~2, but it does not use the complex proxy coefficient directly. Instead, it works with the proxy received-power quantity
\begin{equation}
    \tilde{\rho}_{u,c,n} =
    \frac{P_{\mathrm{AP}}}{P_{\mathrm{UE}}}
    \left|\tilde{h}_{u,c,n}\right|^2,
\end{equation}
where the factor $P_{\mathrm{AP}}/P_{\mathrm{UE}}$ rescales the measured \gls{ue}-to-\gls{ue} proxy power to the distributed-\gls{ap} transmit-power budget. The optimization-3 objective is then
\begin{equation}
    J_w^{(3)}(\mathcal{S}) =
    \frac{1}{|\mathcal{W}_w|\,|\mathcal{U}|}
    \sum_{n \in \mathcal{W}_w}
    \sum_{u \in \mathcal{U}}
    \sum_{c \in \mathcal{S}}
    \tilde{\rho}_{u,c,n}.
\end{equation}
Optimization~3 is therefore the power-only counterpart of optimization~2: it uses the same nearest-user proxy association, but it requires only $|\mathrm{CSI}|^2$ rather than the complex proxy channel.

\paragraph{Greedy subset construction.}
All three optimization rules use the same forward-selection search. Starting from $\mathcal{S}_w^{(0)}=\varnothing$, the controller adds one candidate at a time according to
\begin{equation}
    c_r^\star =
    \arg\max_{c \in \mathcal{C}\setminus \mathcal{S}_w^{(r-1)}}
    J_w^{(q)}\!\left(\mathcal{S}_w^{(r-1)} \cup \{c\}\right),
\end{equation}
\begin{equation}
    q \in \{1,2,3\},
    \qquad
    \mathcal{S}_w^{(r)} =
    \mathcal{S}_w^{(r-1)} \cup \{c_r^\star\},
    \qquad r=1,\ldots,N_{\mathrm{mov}},
\end{equation}
and finally sets $\mathcal{S}_w=\mathcal{S}_w^{(N_{\mathrm{mov}})}$. Consequently, the current implementation is a forward-selection heuristic rather than an exhaustive combinatorial optimizer.

After either optimization rule has chosen the subset, the selected geometric candidate sites are converted into physical relocations of labeled movable \glspl{ap}. Let $\mathbf{p}_m^{(w-1)}$ denote the previous-window position of movable \gls{ap} $m$ and let $\mathbf{q}_j^{(w)}$ denote the coordinates of the candidates in the newly selected subset. The controller solves the linear assignment problem
\begin{equation}
    \pi_w^\star =
    \arg\min_{\pi \in \Pi_{N_{\mathrm{mov}}}}
    \sum_{m=1}^{N_{\mathrm{mov}}}
    \left\|\mathbf{p}_m^{(w-1)}-\mathbf{q}_{\pi(m)}^{(w)}\right\|_2^2,
\end{equation}
where $\Pi_{N_{\mathrm{mov}}}$ is the permutation set on $\{1,\ldots,N_{\mathrm{mov}}\}$. This minimum-displacement matching preserves \gls{ap} identities across windows and yields schedule traces that correspond to physically meaningful AP trajectories instead of arbitrary candidate relabelings.

Once all windows have been simulated, the concatenated movable-strategy trajectory is evaluated using the same trajectory-level summary score as the static baselines,
\begin{equation}
    F = -\alpha P_{\mathrm{out}} + \beta P_{10} + \eta T_{\mathrm{peer}},
\end{equation}
where $P_{\mathrm{out}}$ is the fraction of user snapshots below the configured \gls{sinr} threshold, $P_{10}$ is the trajectory-wide 10th percentile, and $T_{\mathrm{peer}}$ is a peer-weighted mean \gls{sinr} term derived from \gls{ue}-to-\gls{ue} \gls{csi}. The coefficients $(\alpha,\beta,\eta)$ correspond to the configured outage, percentile, and peer tie-break weights. This final score is used for reporting and mode comparison; the online relocation controller itself is driven by one of $J_w^{(1)}(\mathcal{S})$, $J_w^{(2)}(\mathcal{S})$, or $J_w^{(3)}(\mathcal{S})$, depending on the selected movable-strategy variant.
The above formulation defines the relocation mechanisms used throughout the Rabot simulations.

\section{Simulation Setup}

The current paper scaffold is instantiated with the Rabot outdoor scenario defined in \texttt{scenarios/rabot.yaml}. The scene is derived from an \gls{osm} boundary around the KU Leuven Gent Campus Rabot environment. Candidate movable-\gls{ap} positions are generated along building walls at low height, while user mobility is simulated over the walkable outdoor space with pedestrian- and bicycle-like speed profiles.

The simulation pipeline used throughout the manuscript follows the same sequence as the repository workflow: scene construction, candidate generation, user-trajectory generation, path-based \gls{csi} extraction, strategy comparison, and postprocessing into publication-oriented figures. This section should remain descriptive and reproducible; interpretation belongs in the results and discussion sections.

\begin{table}[t]
    \centering
    \caption{Key Rabot scenario parameters used by the current manuscript scaffold.}
    \label{tab:rabot_setup}
    \begin{tabular}{ll}
        \toprule
        Parameter & Value \\
        \midrule
        Carrier frequency & \SI{3.5}{\giga\hertz} \\
        Bandwidth & \SI{100}{\mega\hertz} \\
        User height & \SI{1.5}{\meter} \\
        Number of users & 10 \\
        Simulation horizon & \SI{600}{\second} \\
        Time step & \SI{10}{\second} \\
        Movable \glspl{ap} & 6 \\
        Fixed \glspl{ap} & 0 \\
        Candidate wall height & \SI{1.5}{\meter} \\
        Candidate wall spacing & \SI{18}{\meter} \\
        Minimum candidate spacing & \SI{12}{\meter} \\
        Relocation window & \SI{60}{\second} \\
        Compared strategies & baseline, local heuristic \\
        Exact search & disabled in current Rabot config \\
        Coverage maps & disabled in current Rabot config \\
        \bottomrule
    \end{tabular}
\end{table}

The manuscript should also state the current output contract because the figures in \cref{sec:results_geometry,sec:results_sinr,sec:results_relocation} are directly linked to committed simulation artifacts. In the present workflow, the paper expects scene-layout figures in \texttt{outputs/rabot/} and postprocessed \gls{sinr}/schedule figures in \texttt{outputs/rabot/postprocessing/}.

\gilles{Optional later addition: add one compact pipeline schematic if the venue benefits from a visual overview.}

\section{Scenario Geometry and Mobility Results}
\label{sec:results_geometry}

The first result block should orient the reader in the simulated environment before any performance metric is discussed. The goal is to show where candidate mounting locations exist, how the selected constellation relates to the urban geometry, and how the user trajectories populate the scene over time.

\begin{figure*}[t]
    \centering
    \begin{subfigure}[t]{0.47\textwidth}
        \centering
        \PaperImage[\linewidth]{\resultroot/scene_layout.png}
        \caption{Top-down Rabot scene layout with candidate sites, selected sites, and simulated mobility overlays.}
        \label{fig:scene_layout}
    \end{subfigure}\hfill
    \begin{subfigure}[t]{0.47\textwidth}
        \centering
        \PaperImage[\linewidth]{\resultroot/trajectory_colormap.png}
        \caption{Trajectory colormap used to visualize the temporal spread of user motion through the walkable space.}
        \label{fig:trajectory_colormap}
    \end{subfigure}
    \caption{Geometry-facing result hooks for the Rabot scenario. These figures should remain the first visual anchor for the simulation setup before performance interpretation begins.}
    \label{fig:geometry_pair}
\end{figure*}

The text accompanying \cref{fig:geometry_pair} should stay factual in this draft. It can describe the modeled environment, the spread of the trajectories, and the low-height infrastructure assumption without yet claiming why one region or path is more favorable than another.

\gilles{Later pass: add 2--3 sentences on how the candidate-wall generation and user-motion spread affect the interpretation of later SINR and schedule figures.}
\FloatBarrier

\section{Radio-Performance Distributions}
\label{sec:results_sinr}

This section reports radio-performance statistics for the evaluated placement strategies. The panels summarize the empirical \gls{sinr} distribution across user snapshots and the aggregate \gls{esr} distribution across evaluated time samples.

\begin{figure*}[t]
    \centering

    \ref{sharedlegend}
    \vspace*{1em}
    
    \begin{subfigure}[t]{0.48\textwidth}
        \centering
        \begingroup
        \def\postprocessfigurewidth{\linewidth}
        \pgfplotsset{table/search path={\SINRFigureRoot}}
        \input{\SINRCdfFigurePath}
        \endgroup
        \caption{Empirical \gls{sinr} distribution across \num{7210} user snapshots: \num{10} scheduled users over \num{721} time samples.}
        \label{fig:sinr_cdf}
    \end{subfigure}\hfill
    \begin{subfigure}[t]{0.48\textwidth}
        \centering
        \begingroup
        \def\postprocessfigurewidth{\linewidth}
        \pgfplotsset{table/search path={\SINRFigureRoot}}
        \input{\ESRCdfFigurePath}
        \endgroup
        \caption{\gls{esr} \gls{cdf} across the \num{721} evaluated time samples.}
        \label{fig:esr_cdf}
    \end{subfigure}
    \caption{Radio-performance distributions for the campus trace. \strategy4[title] produces the most favorable \gls{sinr} and \gls{esr} distributions, with \strategy5[title] as the next-best movable strategy.}
    \label{fig:radio_primary_pair}
\end{figure*}

\Cref{fig:radio_primary_pair} establishes a clear ordering. The nearest-user proxy-\gls{esr} controller of \strategy4[title] dominates almost the full empirical \gls{sinr} range. Its mean \gls{sinr} reaches \SI{39.7}{dB}, with a \num{10}th percentile of \SI{24.5}{dB}; \strategy2[title] reaches \SI{21.3}{dB} mean and \SI{7.5}{dB} at the \num{10}th percentile. The outage fraction below \SI{0}{dB} drops from \SI{4.41}{\percent} for \strategy2[title] to \SI{1.87}{\percent} for \strategy4[title]. \strategy5[title] is the second-best radio strategy with \SI{35.7}{dB} mean \gls{sinr} and \SI{16.3}{dB} \num{10}th-percentile \gls{sinr}, but its outage fraction remains essentially identical to \strategy2[title] at about \SI{4.42}{\percent}. Hence, the power-only proxy improves the middle and upper portions of the distribution much more strongly than the weakest-link tail.

\strategy1[title] is weakest throughout, with \SI{9.8}{dB} mean \gls{sinr}, \SI{-22.1}{dB} \num{10}th-percentile \gls{sinr}, and \SI{18.4}{\percent} outage.

The \gls{esr} panel in \cref{fig:radio_primary_pair} follows the same ranking. Mean \gls{esr} reaches \SI{133.1}{bit/s/Hz} for \strategy4[title], \SI{122.8}{bit/s/Hz} for \strategy5[title], \SI{74.9}{bit/s/Hz} for \strategy2[title], and \SI{56.4}{bit/s/Hz} for \strategy1[title]. Relative to \strategy2[title], \strategy4[title] improves the mean \gls{esr} by \SI{77.8}{\percent} and the \num{10}th-percentile \gls{esr} by \SI{88.8}{\percent}. \strategy5[title] still improves the mean \gls{esr} by \SI{64.1}{\percent}, but remains consistently below \strategy4[title].

Together, these figures indicate that the best result in this scenario is not merely ``more movement'' or ``more power'', but the combination of relocation with a phase-aware \gls{ue}-to-\gls{ue} proxy-\gls{csi} objective. \strategy2[title] remains a strong zero-relocation benchmark, whereas \strategy1[title] is clearly mismatched to the spatially spread campus user process.
% \FloatBarrier

\section{Relocation and Practicality Results}
\label{sec:results_relocation}

The relocation figures complement the radio-performance results by quantifying the physical reconfiguration burden associated with movable infrastructure. The scheduling outputs report how often the movable \glspl{ap} change position and how much displacement the resulting schedules induce.

\begin{figure*}[t]
    \centering
    \begin{subfigure}[t]{0.47\textwidth}
        \centering
        \begingroup
        \def\postprocessfigurewidth{\linewidth}
        \pgfplotsset{table/search path={\ScheduleFigureRoot}}
        \input{\ScheduleOverviewFigurePath}
        \endgroup
        \caption{Strategy-level overview of total relocation distance and move fraction.}
        \label{fig:schedule_overview}
    \end{subfigure}\hfill
    \begin{subfigure}[t]{0.47\textwidth}
        \centering
        \begingroup
        \def\postprocessfigurewidth{\linewidth}
        \pgfplotsset{table/search path={\ScheduleFigureRoot}}
        \input{\ScheduleHistogramFigurePath}
        \endgroup
        \caption{Distribution of relocation distances across observed \gls{ap} moves.}
        \label{fig:schedule_histogram}
    \end{subfigure}
    \caption{Relocation-cost figures used to interpret how much physical infrastructure movement is required by the compared placement strategies.}
    \label{fig:schedule_pair}
\end{figure*}

Together, the panels in \cref{fig:schedule_pair} characterize the relocation activity associated with the evaluated placement strategies through aggregate movement summaries and relocation-distance distributions.
\FloatBarrier

\section{Discussion and Current Limits}

This draft deliberately stops short of a full interpretation layer. Still, the manuscript should state the boundaries of the current evidence clearly so that later quantitative claims remain anchored to what the simulation pipeline actually models.

The main limitations that should remain visible in the first paper version are the following:
\begin{itemize}
    \item Only a single outdoor scenario is wired into the current manuscript scaffold.
    \item The current Rabot configuration disables coverage-map exports, so the paper centers on trajectory-level \gls{sinr} and schedule metrics.
    \item The current checked scenario also disables capped exact search, so the default draft compares the baseline and the local \gls{csi}-driven heuristic unless an updated output set is committed.
    \item Citizen participation is abstracted as movable candidate positions and scheduler actions rather than being validated experimentally.
    \item The present text does not yet discuss sensitivity to seeds, scenario perturbations, or broader deployment costs.
\end{itemize}

The next analysis pass should therefore answer at least three questions: how stable the reported gains are across seeds and scenarios, how strongly the conclusions depend on the chosen \gls{sinr} objective, and how the relocation metrics should be interpreted in a real co-created deployment process.

\gilles{Later pass: split this section into limitations and future work if the venue expects a more formal discussion structure.}

\section{Conclusion}

This article presented a simulation-based study of low-height distributed-\gls{mimo} access networks with repositionable infrastructure in an outdoor urban scenario. The paper combined a system model, three relocation algorithms, a reproducible Rabot simulation workflow, and postprocessed geometry, \gls{sinr}, \gls{esr}, and relocation figures derived from exported artifacts. The three movable strategies span a historical local-\gls{ap}-to-\gls{ue} \gls{csi} controller, a nearest-user \gls{ue}-to-\gls{ue} proxy-\gls{csi} controller with a proxy-\gls{esr} objective, and a nearest-user proxy-power controller based on $|\mathrm{CSI}|^2$.

Within its current scope, the framework supports systematic comparison of a central rooftop proxy, a fixed distributed deployment, and three relocation-capable distributed deployments under a common channel model and mobility process. Natural next steps are multi-scenario studies, coverage-map analysis, sensitivity analysis for the proxy-based relocation rules, and experimental validation in the broader Cocoon program.


% Optional appendix hook for future coverage-map figures.
% \appendices
% \section{Optional Coverage Maps}
% \PaperSubfigurePair[t][0.48\textwidth]
%   {\FixedCoverageFigurePath}{Baseline coverage map.}{fig:coverage_baseline}
%   {\CoverageFigurePath}{Best-strategy coverage map.}{fig:coverage_best}
%   {Optional coverage-map comparison. Enable this appendix once coverage exports are part of the paper narrative.}{fig:coverage_pair}

\section*{Acknowledgment}
\gilles{Add funding statement, collaborators, and infrastructure acknowledgments here.}

\printbibliography

\end{document}
